\section{Lagrange multipliers}

Lagrange multipliers find extrema of a function subject to equality constraints. The constraints restrict the search from the ambient space to a curve, surface, or higher-dimensional constraint set. Let $f(\mathbf{x})$ be the objective and let
\begin{equation}
c_a(\mathbf{x})=0,
\qquad a=1,\ldots,k,
\end{equation}
be independent constraints. Introduce the augmented function
\begin{equation}
\mathcal F(\mathbf{x},\lambda_a)=f(\mathbf{x})-\lambda_a c_a(\mathbf{x}),
\end{equation}
where repeated $a$ is summed. Treating both $\mathbf{x}$ and the multipliers as independent variables gives the stationary equations
\begin{equation}
\frac{\partial f}{\partial x_i}
-\lambda_a\frac{\partial c_a}{\partial x_i}=0,
\qquad
c_a(\mathbf{x})=0.
\end{equation}

The multiplier values are not chosen in advance; they are determined along with the extremizing point. Geometrically, allowed displacements satisfy
\begin{equation}
\mathrm{d}c_a=\nabla c_a\cdot\mathrm{d}\mathbf{x}=0.
\end{equation}
At a constrained extremum, $\nabla f$ is orthogonal to every allowed displacement and therefore lies in the span of the constraint normals:
\begin{equation}
\nabla f=\lambda_a\nabla c_a.
\end{equation}

\begin{example}[Closest point on a plane]
Find the point on $x+y+z=1$ closest to the origin. Minimize
\begin{equation}
f=x^2+y^2+z^2
\end{equation}
subject to $c=x+y+z-1=0$. The multiplier equations are
\begin{equation}
2x-\lambda=0,
\qquad 2y-\lambda=0,
\qquad 2z-\lambda=0,
\qquad x+y+z=1.
\end{equation}
Thus $x=y=z$, and the closest point is $(1/3,1/3,1/3)$.
\end{example}

This geometric statement explains why the method works: moving along the constraint surface has no first-order effect on $f$, even though $\nabla f$ need not vanish in the full ambient space. The same method applies to several constraints, including extrema on a sphere, distance to a surface, and optimization with normalization conditions. The Hessian of the objective, $H_{ij}=\partial_i\partial_jf$, determines the local type of an unconstrained nondegenerate critical point; constrained classification requires the Hessian restricted to allowed directions.